\documentclass[10pt,letterpaper]{article}
\usepackage[textwidth=5.5in,textheight=9in,centering]{geometry}
\usepackage{times,amsmath,amssymb,amsthm,graphicx,microtype,flafter,placeins}
\usepackage[colorlinks=true,linkcolor=blue,urlcolor=blue]{hyperref}
\usepackage{titlesec}
\titleformat{\subsection}{\normalfont\normalsize\scshape}{\thesubsection}{1em}{}
\setlength{\parindent}{0pt}
\setlength{\parskip}{5pt plus 1pt minus 1pt}
\setlength{\textfloatsep}{10pt plus 2pt minus 2pt}
\setlength{\intextsep}{10pt plus 2pt minus 2pt}
\widowpenalty=10000
\clubpenalty=10000
\newtheorem{theorem}{Theorem}[section]
\newtheorem{proposition}{Proposition}[section]
\hypersetup{pdftitle={Why joint training acquires scale slowly},pdfauthor={}}
\begin{document}
% Match the spectrum note's proposed insertion point. The included files
% carry no manuscript counters or manual figure/equation numbers.
\setcounter{section}{3}
\setcounter{subsection}{4}
\setcounter{theorem}{2}
\setcounter{figure}{3}
\input{docs/d34_scale_acquisition_paper_main.tex}
\FloatBarrier
\clearpage
\appendix
\setcounter{figure}{0}
\renewcommand{\thefigure}{S\arabic{figure}}
\numberwithin{equation}{section}
\input{docs/d34_scale_acquisition_paper_appendix.tex}
\end{document}
